% Generated draft: requires mathematical and source review.
\documentclass[11pt]{article}
\usepackage[margin=25mm]{geometry}
\usepackage{fontspec}
\setmainfont{Latin Modern Roman}
\usepackage{amsmath,amssymb,mathtools,unicode-math}
\setmathfont{Latin Modern Math}
\usepackage{enumitem,needspace,xurl,hyperref}
\setlength{\parindent}{0pt}
\setlength{\parskip}{5pt}
\setlength{\emergencystretch}{3em}
\newcommand{\N}{\mathbb{N}}
\newcommand{\Z}{\mathbb{Z}}
\newcommand{\Q}{\mathbb{Q}}
\newcommand{\R}{\mathbb{R}}
\newcommand{\C}{\mathbb{C}}
\newcommand{\F}{\mathbb{F}}
\newcommand{\A}{\mathbb{A}}
\newcommand{\PP}{\mathbb{P}}
\newcommand{\E}{\mathbb{E}}
\newcommand{\eps}{\varepsilon}
\newcommand{\dd}{\,\mathrm{d}}
\begin{document}
\section*{UnsolvedMath ID: 701700}
\textbf{Source problem code:} catalogue entry
\subsection{Definitions and mathematical statement}
Let \(G\) be a group. It is finitely generated if \(G=\langle S\rangle\) for some finite subset \(S\subseteq G\). It is residually finite if, for every \(g\in G\setminus\{1_G\}\), there exist a finite group \(H\) and a homomorphism \(\varphi:G\to H\) such that \(\varphi(g)\ne 1_H\).

For an integer \(r\ge 1\), let \(F_r=\langle x_1,\ldots,x_r\rangle\) be the free group on \(r\) generators. The group \(G\) satisfies a nontrivial group law if there exist \(r\ge 1\) and \(w\in F_r\setminus\{1_{F_r}\}\) such that
\[
\forall (g_1,\ldots,g_r)\in G^r,\qquad w(g_1,\ldots,g_r)=1_G.
\]
Here \(w(g_1,\ldots,g_r)\) denotes evaluation under the homomorphism \(F_r\to G\) sending each \(x_i\) to \(g_i\).

Regard groups as discrete. A group \(G\) is amenable if there exists a finitely additive probability measure \(\mu:\mathcal P(G)\to[0,1]\) on all subsets of \(G\) such that \(\mu(gA)=\mu(A)\) for every \(g\in G\) and \(A\subseteq G\). Thus \(\mu(G)=1\), \(\mu(\varnothing)=0\), and \(\mu(A\cup B)=\mu(A)+\mu(B)\) for disjoint subsets \(A,B\subseteq G\). The class of elementary amenable groups is the smallest class containing all finite groups and all abelian groups and closed under taking subgroups, quotients, group extensions, and directed unions.

\textbf{Existence question with an unresolved parenthetical.} The supplied statement asks whether there exists a finitely generated residually finite group which satisfies a nontrivial group law and is ``not (elementary) amenable.'' This wording admits the following two targets:
\begin{enumerate}
\item Does there exist a finitely generated residually finite group \(G\) satisfying a nontrivial group law such that \(G\) is not amenable?
\item Does there exist a finitely generated residually finite group \(G\) satisfying a nontrivial group law such that \(G\) is not elementary amenable?
\end{enumerate}
The supplied record does not specify whether both targets are intended or whether the parenthetical selects one of them. Every elementary amenable group is amenable, so an example for the first target would also satisfy the second; the two qualifications must nevertheless be distinguished. No equivalence of the targets is asserted.
\subsection{Short English statement}
Does there exist a finitely generated residually finite group satisfying a nontrivial group identity that is nonamenable, or, under the alternative reading of the supplied parenthetical, not elementary amenable? The intended role of ``(elementary)'' requires clarification. This is a statement-only draft requiring human mathematical review.
\subsection{Sources}
\begin{itemize}
\item Local UnsolvedMath registry, permanent ID \(701700\), exact source code \texttt{catalogue entry}, title \emph{Kourovka Notebook Problem 20.42}. The supplied record labels the problem \texttt{open}; this is a source label, not an independently verified status assessment.
\item Supplied source attribution: \emph{The Kourovka Notebook (Unsolved Problems in Group Theory)}, 20th issue (2022), original problem number 20.42; also identified with the 21st issue edition (2026), \texttt{arXiv:1401.0300v48}. The collection metadata names E. I. Khukhro and V. D. Mazurov as editors. These bibliographic details are supplied and have not been independently checked.
\item The background and proposer metadata attribute the problem to S. Kionke, with proposed year 2022. The statement itself begins with the parenthetical attribution ``Y. Cornulier, A. Mann, A. Thom.'' Both attributions are retained; their respective roles are unspecified.
\end{itemize}
\end{document}
